\section{Preliminaries}

For a metric space $X$, let $c_2(X)$ denote its least bilipschitz distortion
into Hilbert space, with $c_2(X)=\infty$ if no such embedding exists.  We use the
standard finite-determination principle
\begin{equation}\label{eq:finite-determination}
        c_2(X)=\sup\{c_2(F):F\subset X,\ |F|<\infty\},
\end{equation}
which is \cite[Proposition~31]{CahillIversonMixon} and follows from the usual
Hilbert-ultraproduct argument.

For $n\ge1$, put
\[
        H_n=\ell^2(\Z^n),
        \qquad
        \Sigma_n=\Sphere(H_n)/\Z^n,
\]
where $\Z^n$ acts by translations
\[
        (\tau_a f)(k)=f(k-a),\qquad a,k\in\Z^n.
\]
Unless stated otherwise, $\Sigma_n$ carries the chordal quotient metric
\[
        d_{\Sigma_n}([f],[g])=\inf_{a\in\Z^n}\norm{f-\tau_a g}_2.
\]
The same argument applies in $H_n$: if $\tau_{a_j}f\to g$ and, after passing
to a subsequence, $a_j\to\infty$ in $\Z^n$, then $\tau_{a_j}f\rightharpoonup0$,
contradicting $\norm{g}=1$; otherwise the sequence $(a_j)$ has a constant
subsequence.  Thus translation orbits on $\Sphere(H_n)$ are closed.
For unit vectors this is equivalently
\begin{equation}\label{eq:sphere-correlation}
        d_{\Sigma_n}([f],[g])^2
        =2-2\sup_{a\in\Z^n}\ip{f}{\tau_a g}.
\end{equation}

The only external input is the flat-torus lower bound of Khot--Naor: there
is a universal constant $c>0$ and, for arbitrarily large $n$, a lattice
$L_n\subset\R^n$ such that
\begin{equation}\label{eq:KN}
        c_2(\R^n/L_n)\ge c\sqrt n .
\end{equation}
One may obtain this directly from \cite[Theorem~1.1 and the paragraph following
it]{HavivRegev}, where the theorem is attributed to Khot--Naor.  In the original
paper, this is the flat-torus consequence of \cite[Theorem~5.1 and
Corollary~5.2]{KhotNaor}.

Given $A\in GL_n(\R)$, write $T_A=\R^n/\Z^n$ with the flat quotient metric
\begin{equation}\label{eq:rhoA}
        \rho_A(u,v)=\min_{m\in\Z^n}\norm{A(u-v-m)}_2.
\end{equation}
Then $T_A$ is isometric to $\R^n/A\Z^n$: if $L=A\Z^n$, the map
$x+L\mapsto A^{-1}x+\Z^n$ identifies $\R^n/L$ with $T_A$.  Thus the hard
flat tori supplied by \eqref{eq:KN} may be regarded as tori with the unit
lattice; the matrix $A$ simply records the flat metric.

We shall use the following standard form of Poisson summation.

\begin{lemma}[Uniform summation estimate]\label{lem:summation}
Let $G\in\mathcal S(\R^n)$.  For every $N\ge1$ there is a constant
$C_{G,N}<\infty$ such that, for every $R\ge1$ and every $c\in\R^n$,
\[
\left|R^{-n}\sum_{k\in\Z^n}G\left(\frac{k-c}{R}\right)
      -\int_{\R^n}G(x)\,dx\right|
\le C_{G,N}R^{-N}.
\]
\end{lemma}

\begin{proof}
With the convention
$\widehat G(\xi)=\int_{\R^n}G(x)e^{-2\pi i\ip{x}{\xi}}\,dx$, Poisson summation
gives
\[
        R^{-n}\sum_{k\in\Z^n}G\left(\frac{k-c}{R}\right)
        =\sum_{\ell\in\Z^n}e^{-2\pi i\ip{\ell}{c}}\widehat G(R\ell).
\]
The term $\ell=0$ is $\int_{\R^n}G$.  Since $\widehat G$ is Schwartz, choose
$M>N+n$ and bound the remaining terms by
\[
        \sum_{\ell\ne0}|\widehat G(R\ell)|
        \le C_MR^{-M}\sum_{\ell\ne0}\norm{\ell}^{-M}
        \le C_{G,N}R^{-N},
\]
uniformly in $c$.
\end{proof}
